%%%%%%%%%%%%%%%%%%%%%%%%%%%%%%%%%%%%%%%%%%%%%%%%%%%%%%%%%%%%%%%%%%%%%%%%%%%%%%%%%%
\begin{frame}[fragile]\frametitle{}
\begin{center}
{\Large Getting Started}
\end{center}
\end{frame}

%%%%%%%%%%%%%%%%%%%%%%%%%%%%%%%%%%%%%%%%%%%%%%%%%%%%%%%%%%%
\begin{frame}[fragile]\frametitle{Learning Roadmap for Freshers}

\begin{enumerate}
  \item \textbf{Step 1: Foundations (Weeks 1--2)}
    \begin{itemize}
      \item Linear algebra refresher: matrices, eigenvalues, dot products.
      \item Graph theory basics: nodes, edges, adjacency matrix, graph Laplacian.
      \item PyTorch basics if you don't already know it.
    \end{itemize}
  \item \textbf{Step 2: First GNN (Weeks 3--4)}
    \begin{itemize}
      \item Work through the PyTorch Geometric (PyG) tutorials.
      \item Implement GCN on Cora (node classification), you saw the code today.
      \item Read the original GCN paper (Kipf \& Welling, 2017), only 9 pages.
    \end{itemize}
  \item \textbf{Step 3: Go Deeper (Months 2--3)}
    \begin{itemize}
      \item Stanford CS224W (free online): Machine Learning with Graphs.
      \item Read the Bronstein et al.\ (2021) GDL survey paper.
    \end{itemize}
  \item \textbf{Step 4: Specialise}
    \begin{itemize}
      \item Pick one domain: molecules, social networks, 3D vision, or NLP graphs.
      \item Reproduce one paper from that domain end-to-end.
    \end{itemize}
\end{enumerate}

\end{frame}

%%%%%%%%%%%%%%%%%%%%%%%%%%%%%%%%%%%%%%%%%%%%%%%%%%%%%%%%%%%
\begin{frame}[fragile]\frametitle{How to Choose Your Architecture}

\begin{enumerate}
  \item \textbf{Identify your data geometry}
    \begin{itemize}
      \item What symmetries does your data exhibit?
      \item What transformations should leave the output unchanged?
      \item What is the natural neighbourhood structure?
    \end{itemize}
  \item \textbf{Choose your geometric domain}
    \begin{itemize}
      \item Regular structure? $\rightarrow$ CNN (Grid)
      \item Global symmetries (sphere, rotation)? $\rightarrow$ Group-equivariant network
      \item Irregular connectivity? $\rightarrow$ GNN (Graph)
      \item Curved 3D surface? $\rightarrow$ Mesh network (Geodesic)
    \end{itemize}
  \item \textbf{Design the architecture using the blueprint}
    \begin{itemize}
      \item Local equivariant layers $\rightarrow$ Non-linearity $\rightarrow$ Local pooling $\rightarrow$ Global pooling
    \end{itemize}
\end{enumerate}

\end{frame}

%%%%%%%%%%%%%%%%%%%%%%%%%%%%%%%%%%%%%%%%%%%%%%%%%%%%%%%%%%%
\begin{frame}[fragile]\frametitle{Key Libraries and Tools}

\begin{columns}
\begin{column}{0.5\textwidth}
\textbf{Graph Neural Networks:}
\begin{itemize}
  \item \textbf{PyTorch Geometric (PyG)}: Most complete, production-ready. Start here.
  \item \textbf{DGL}: Deep Graph Library (Facebook). Fast and flexible.
  % \item \textbf{Spektral}: Keras/TF-based graph learning.
\end{itemize}

\textbf{Manifold Learning:}
\begin{itemize}
  \item \textbf{Geomstats}: Computation on nonlinear manifolds.
  \item \textbf{DiffusionNet}: Learning on meshes.
  \item \textbf{PyManOpt}: Optimisation on manifolds.
\end{itemize}
\end{column}
\begin{column}{0.5\textwidth}
\textbf{Group-Equivariant Networks:}
\begin{itemize}
  \item \textbf{e3nn}: Euclidean neural networks (SE(3)).
  \item \textbf{escnn}: Equivariant steerable CNNs.
\end{itemize}

\textbf{Topology:}
\begin{itemize}
  \item \textbf{GUDHI}: Topological data analysis.
  \item \textbf{scikit-tda}: TDA for machine learning.
\end{itemize}

\textbf{Symbolic Geometry:}
\begin{itemize}
  \item \textbf{SymPy}: Symbolic math.
  \item \textbf{diffgeom}: Differential geometry (PyPi).
\end{itemize}
\end{column}
\end{columns}

\begin{alertblock}{Start Here}
\centering PyTorch Geometric provides the most comprehensive and well-documented entry point for GDL.
\end{alertblock}

\end{frame}
